\documentclass[11pt,letterpaper]{article}

\usepackage[T1]{fontenc}
\usepackage{amsmath,amssymb}
\usepackage{tikz}
\usepackage{pgfplots}
\pgfplotsset{compat=1.18}
\usepackage{algorithm}
\usepackage{algpseudocode}
\usepackage{booktabs}
\usepackage{enumitem}
\usepackage[hidelinks]{hyperref}

\newcommand{\N}{\mathbb{N}}
\newcommand{\Z}{\mathbb{Z}}
\newcommand{\Q}{\mathbb{Q}}
\newcommand{\R}{\mathbb{R}}
\newcommand{\C}{\mathbb{C}}
\newcommand{\Map}[3]{#1 \colon #2 \to #3}
\newcommand{\Mod}[1]{#1\text{-}\mathrm{Mod}}
\newcommand{\GroupHom}{\mathrm{Hom}_{\mathbf{Grp}}}
\newcommand{\RingHom}{\mathrm{Hom}_{\mathbf{Ring}}}

\newtheorem{theorem}{Theorem}
\newtheorem{lemma}[theorem]{Lemma}
\newtheorem{corollary}[theorem]{Corollary}
\newtheorem{definition}[theorem]{Definition}
\newtheorem{remark}[theorem]{Remark}
\newenvironment{proof}{\noindent\textit{Proof.}}{\hfill\(\square\)}

\title{Five displayed solutions of \(x^{2}+7=2^{n}\)}
\author{}
\date{\today}

\begin{document}
\maketitle

\begin{abstract}
The five pairs \((1,3)\), \((3,4)\), \((5,5)\), \((11,7)\), and \((181,15)\) satisfy \(x^{2}+7=2^{n}\). Each proof evaluates one square and one power of two. The same arithmetic does not show that every positive solution is one of these pairs.
\end{abstract}

\section{The equation}

Let \(\N\) be the nonnegative integers. For \(x,n\in\N\), write \(E(x,n)\) for the equation \(x^{2}+7=2^{n}\). With the exponent fixed, this is the statement that the natural number \(x\) is a root of the polynomial \(X^{2}+7-2^{n}\).

\begin{definition}
The five displayed pairs are $(1,3)$, $(3,4)$, $(5,5)$, $(11,7)$, and $(181,15)$.
\end{definition}

\section{Powers of two}

\begin{lemma}
\label{lem:powers}
\(2^{3}=8\), \(2^{4}=16\), \(2^{5}=32\), \(2^{6}=64\), \(2^{7}=128\), and \(2^{15}=32768\).
\end{lemma}

\begin{proof}
The rule \(2^{k+1}=2^{k}\cdot 2\), starting from \(2^{1}=2\), gives \(2^{2}=4\), \(2^{3}=8\), \(2^{4}=16\), \(2^{5}=32\), \(2^{6}=64\), and \(2^{7}=128\). The exponent splits \(10=5+5\) and \(15=10+5\), so
\[
2^{10}=2^{5}\cdot 2^{5}=32\cdot 32=1024
\]
and
\[
2^{15}=2^{10}\cdot 2^{5}=1024\cdot 32=32768.
\]
\end{proof}

\section{Squares}

\begin{lemma}
\label{lem:squares}
\(1^{2}=1\), \(3^{2}=9\), \(5^{2}=25\), \(11^{2}=121\), and \(181^{2}=32761\).
\end{lemma}

\begin{proof}
The first three products are \(1\cdot 1=1\), \(3\cdot 3=9\), and \(5\cdot 5=25\). For the fourth, \(11=10+1\), so
\begin{align*}
11\cdot 11
&=(10+1)(10+1)\\
&=(10+1)\cdot 10+(10+1)\cdot 1\\
&=(10\cdot 10+1\cdot 10)+(10\cdot 1+1\cdot 1)\\
&=121.
\end{align*}
For the fifth, \(180=18\cdot 10\) and \(18\cdot 18=324\), hence
\begin{align*}
180\cdot 180
&=(18\cdot 10)(18\cdot 10)\\
&=18\cdot(10\cdot(18\cdot 10))\\
&=18\cdot(18\cdot(10\cdot 10))\\
&=(18\cdot 18)\cdot(10\cdot 10)\\
&=324\cdot 100\\
&=32400.
\end{align*}
Then \(181=180+1\), so
\begin{align*}
181\cdot 181
&=(180+1)(180+1)\\
&=(180+1)\cdot 180+(180+1)\cdot 1\\
&=(180\cdot 180+1\cdot 180)+(180\cdot 1+1\cdot 1)\\
&=32761.
\end{align*}
\end{proof}

\section{The five pairs}

\begin{theorem}
\label{thm:displayed}
Each displayed pair satisfies \(E(x,n)\):
\begin{align*}
1^{2}+7&=2^{3},\\
3^{2}+7&=2^{4},\\
5^{2}+7&=2^{5},\\
11^{2}+7&=2^{7},\\
181^{2}+7&=2^{15}.
\end{align*}
\end{theorem}

\begin{proof}
Lemma~\ref{lem:squares} and Lemma~\ref{lem:powers} give
\(1+7=8\), \(9+7=16\), \(25+7=32\), \(121+7=128\), and \(32761+7=32768\).
These are the five powers \(2^{3}\), \(2^{4}\), \(2^{5}\), \(2^{7}\), and \(2^{15}\).
\end{proof}

\begin{remark}
Theorem~\ref{thm:displayed} is an evaluation at five pairs.
It does not assert that a positive solution of \(x^{2}+7=2^{n}\) equals one of those pairs.
It does not classify solutions of \(x^{2}+D=2^{n}\) for a general positive integer \(D\).
It does not produce a cycle, a congruence, or a positivity window.
\end{remark}

\end{document}
